\section{Setting and the Single-Visit Constraint}
\label{sec:setting}

A click $i$ occurs at time $C_i$ with features $x_i$ and, if it converts, does so after a delay
$D_i>0$, with $D_i=\infty$ when it never converts. The attribution window $v$ belongs to the
objective and fixes the target $\Yv_i=\mathbf{1}[D_i\le v]$, knowable only at $C_i+v$. The join delay $o$ belongs to the pipeline and fixes when the
record is written, so the label available at that moment is
\begin{equation}
  \Yo_i \;=\; \mathbf{1}\!\left[D_i \le o\right] \;\le\; \Yv_i .
  \label{eq:labels}
\end{equation}
Setting $o=v$ reproduces the conventional per-objective pipeline. Setting $o\ll v$ is what we
want, and the inequality in Eq.~\eqref{eq:labels} is its price. We write
$F(o)=P(D\le o\mid D\le v)$ for the fraction of the objective's conversions already visible when
the record is written.

Training proceeds in hourly rounds. Each round consumes the records released since the previous
cut-off, then scores every click of the following hour exactly once against the matured target
$\Yv$. Every method is therefore judged on the business label whatever it trained on, and
updating before predicting keeps the evaluation causal.

\paragraph{Two rules.}
A deployable pipeline must satisfy two rules.

\textbf{Rule 1, causality.} A record's label must be computable at the moment the record is
written, and must not change afterwards. Assigning $\Yv_i$ to a record written at $C_i+o$ uses
an outcome that has not yet occurred. Such an arm is an upper reference, not a system, and we
report it under the name \textsc{Oracle}.

\textbf{Rule 2, single-visit.} Let $n_i$ be the number of times click $i$ enters the served
model's loss over its lifetime. Rule~2 requires $\max_i n_i\le 1$. Without it the question is
empty, since any join delay can be shortened by writing a record early and writing it again
later.

We take the maximum over clicks rather than the mean, because it is the existence of a second
write that obliges the pipeline to support revision, and the rule governs the served stream
only: charging training-time auxiliaries would eliminate \textsc{Fsiw}, \textsc{Ulc} and our
own arms alike.

\paragraph{Which published methods comply.}
We reproduced \nScreened{} published methods from scratch under a single protocol and counted
the records each one's served model consumes per distinct click. \nFail{} write more than
one.\footnote{Two measurement traps both understate duplication. The denominator must be
distinct clicks touched, since dividing records consumed by records released conflates
duplication with variation in click volume. And a scoring window shorter than a method's revisit
horizon cannot observe the revisit, so we score only clicks whose entire horizon lies inside the
stream.} At $v=1$\,d the counts are \textsc{FTP}~\cite{li2021follow}~4.00, \textsc{MISS}~\cite{liu2024online}~3.05, \textsc{PI}~\cite{zhang2025personalized}~2.00, \textsc{DDFM}~\cite{dai2023dually}~1.89, \textsc{DFM}~\cite{chapelle2014modeling}~1.33, \textsc{DEFUSE}~\cite{chen2022asymptotically}~1.03, \textsc{ES{-}DFM}~\cite{yang2021capturing}~1.03
. The highest, \maxVisitsArm{},
reaches \maxVisits{} writes for a single click, its mechanism being an update applied whenever a
label reverses. The plain re-ingestion
baseline~\cite{ktena2019addressing,gu2021real} measures \vanillaDup{} and fails the same
criterion. Duplication is a property of emission logic rather than of the model, and the audit
returns identical counts at the smallest and largest served capacities.

Only \passList{} satisfy Rule~2. Of these, \textsc{Fsiw} corrects the loss and \textsc{Twice}
adds a prediction head on the conversion clock, both leaving the target itself untouched.
\textsc{Ulc} alone changes the target, which makes it our closest comparator. Section~\ref{sec:results}
reports how all of them perform, rather than only whether they comply.
